\documentclass[12pt]{article}

\usepackage[utf8]{inputenc}
\usepackage[T1]{fontenc}
\usepackage[spanish]{babel}
\usepackage{graphicx}
\usepackage{geometry}
\usepackage{fancyhdr}
\usepackage{enumitem}
\usepackage{amsmath,amssymb}
\usepackage{lmodern}
\usepackage{multicol}
\usepackage{xcolor}

\geometry{letterpaper, margin=2cm}

\begin{document}
	
	\vspace*{1em}
	
	\begin{center}
		{\Large \textbf{Soluciones Guía 03}}\\[0.3em]
		{\small Unidades de medida - Expresiones algebraicas}
	\end{center}
	
	% =========================
	% SOLUCIONES I
	% =========================
	\section*{I. Soluciones: Conversiones de unidades}
	
	\begin{enumerate}[itemsep=0.45cm]
		\item \textbf{Convertir $5000\ \text{mL}$ a litros}
		\[
		1000\ \text{mL}=1\ \text{L}
		\qquad\Rightarrow\qquad
		5000\ \text{mL}=\frac{5000}{1000}=5\ \text{L}
		\]
		\textbf{Respuesta:} $5\ \text{L}$.
		
		\item \textbf{Convertir $3.5\ \text{kg}$ a gramos}
		\[
		1\ \text{kg}=1000\ \text{g}
		\qquad\Rightarrow\qquad
		3.5\ \text{kg}=3.5\cdot1000=3500\ \text{g}
		\]
		\textbf{Respuesta:} $3500\ \text{g}$.
		
		\item \textbf{Convertir $750\ \text{g}$ a kilogramos}
		\[
		1000\ \text{g}=1\ \text{kg}
		\qquad\Rightarrow\qquad
		750\ \text{g}=\frac{750}{1000}=0.75\ \text{kg}
		\]
		\textbf{Respuesta:} $0.75\ \text{kg}$.
		
		\item \textbf{Convertir $2.5\ \text{L}$ a mililitros}
		\[
		1\ \text{L}=1000\ \text{mL}
		\qquad\Rightarrow\qquad
		2.5\ \text{L}=2.5\cdot1000=2500\ \text{mL}
		\]
		\textbf{Respuesta:} $2500\ \text{mL}$.
		
		\item \textbf{Convertir $4500\ \text{m}$ a kilómetros}
		\[
		1000\ \text{m}=1\ \text{km}
		\qquad\Rightarrow\qquad
		4500\ \text{m}=\frac{4500}{1000}=4.5\ \text{km}
		\]
		\textbf{Respuesta:} $4.5\ \text{km}$.
		
		\item \textbf{Convertir $2.4\ \text{km}$ a metros}
		\[
		1\ \text{km}=1000\ \text{m}
		\qquad\Rightarrow\qquad
		2.4\ \text{km}=2.4\cdot1000=2400\ \text{m}
		\]
		\textbf{Respuesta:} $2400\ \text{m}$.
		
		\item \textbf{Convertir $180\ \text{min}$ a horas}
		\[
		60\ \text{min}=1\ \text{h}
		\qquad\Rightarrow\qquad
		180\ \text{min}=\frac{180}{60}=3\ \text{h}
		\]
		\textbf{Respuesta:} $3\ \text{h}$.
		
		\item \textbf{Convertir $2.5\ \text{h}$ a minutos}
		\[
		1\ \text{h}=60\ \text{min}
		\qquad\Rightarrow\qquad
		2.5\ \text{h}=2.5\cdot 60=150\ \text{min}
		\]
		\textbf{Respuesta:} $150\ \text{min}$.
		
		\item \textbf{Convertir $3500\ \text{cm}$ a metros}
		\[
		100\ \text{cm}=1\ \text{m}
		\qquad\Rightarrow\qquad
		3500\ \text{cm}=\frac{3500}{100}=35\ \text{m}
		\]
		\textbf{Respuesta:} $35\ \text{m}$.
		
		\item \textbf{Convertir $2500\ \text{kg}$ a toneladas}
		\[
		1000\ \text{kg}=1\ \text{t}
		\qquad\Rightarrow\qquad
		2500\ \text{kg}=\frac{2500}{1000}=2.5\ \text{t}
		\]
		\textbf{Respuesta:} $2.5\ \text{t}$.
		
		\item \textbf{Convertir $900\ \text{mL}$ a litros}
		\[
		1000\ \text{mL}=1\ \text{L}
		\qquad\Rightarrow\qquad
		900\ \text{mL}=\frac{900}{1000}=0.9\ \text{L}
		\]
		\textbf{Respuesta:} $0.9\ \text{L}$.
		
		\item \textbf{Convertir $80000\ \text{cm}^2$ a metros cuadrados}
		\[
		1\ \text{m}^2=(100\ \text{cm})^2=10000\ \text{cm}^2
		\]
		\[
		80000\ \text{cm}^2=\frac{80000}{10000}=8\ \text{m}^2
		\]
		\textbf{Respuesta:} $8\ \text{m}^2$.
		
		\item \textbf{Convertir $3\ \text{ha}$ a metros cuadrados}
		\[
		1\ \text{ha}=10000\ \text{m}^2
		\qquad\Rightarrow\qquad
		3\ \text{ha}=3\cdot10000=30000\ \text{m}^2
		\]
		\textbf{Respuesta:} $30000\ \text{m}^2$.
		
		\item \textbf{Convertir $15000\ \text{m}^2$ a hectáreas}
		\[
		10000\ \text{m}^2=1\ \text{ha}
		\qquad\Rightarrow\qquad
		15000\ \text{m}^2=\frac{15000}{10000}=1.5\ \text{ha}
		\]
		\textbf{Respuesta:} $1.5\ \text{ha}$.
		
		\item \textbf{Convertir $2\ \text{m}^3$ a litros}
		\[
		1\ \text{m}^3=1000\ \text{L}
		\qquad\Rightarrow\qquad
		2\ \text{m}^3=2\cdot1000=2000\ \text{L}
		\]
		\textbf{Respuesta:} $2000\ \text{L}$.
		
		\item \textbf{Convertir $5000\ \text{L}$ a metros cúbicos}
		\[
		1000\ \text{L}=1\ \text{m}^3
		\qquad\Rightarrow\qquad
		5000\ \text{L}=\frac{5000}{1000}=5\ \text{m}^3
		\]
		\textbf{Respuesta:} $5\ \text{m}^3$.
	\end{enumerate}
	
	\newpage
	
	% =========================
	% SOLUCIONES III
	% =========================
	\section*{III. Soluciones: Problemas contextualizados}
	
	\begin{enumerate}[itemsep=0.7cm]
		\item \textbf{Sistema de riego}
		
		La tasa es $1200\ \text{mL/min}$.
		
		Primero convertimos a litros:
		\[
		1200\ \text{mL/min}=1.2\ \text{L/min}
		\]
		
		Ahora convertimos a litros por hora:
		\[
		1.2\ \text{L/min}\times 60\ \text{min/h}=72\ \text{L/h}
		\]
		
		\textbf{Respuesta:} $72\ \text{L/h}$.
		
		\item \textbf{Estanque agrícola}
		
		El volumen del prisma rectangular es:
		\[
		V=largo \cdot ancho \cdot profundidad
		\]
		\[
		V=4\cdot2\cdot1.5=12\ \text{m}^3
		\]
		
		Como:
		\[
		1\ \text{m}^3=1000\ \text{L}
		\]
		entonces:
		\[
		12\ \text{m}^3=12000\ \text{L}
		\]
		
		\textbf{Respuesta:} se necesitan $12000\ \text{L}$ de agua.
		
		\item \textbf{Masa total utilizada}
		
		Fertilizante:
		\[
		3\cdot25=75\ \text{kg}
		\]
		
		Compost:
		\[
		5\cdot18=90\ \text{kg}
		\]
		
		Masa total:
		\[
		75+90=165\ \text{kg}
		\]
		
		Ahora convertimos a toneladas:
		\[
		165\ \text{kg}=\frac{165}{1000}=0.165\ \text{t}
		\]
		
		\textbf{Respuesta:} $0.165\ \text{t}$.
		
		\item \textbf{Velocidad del tractor}
		
		\begin{itemize}
			\item \textbf{En km/h}
			
			Recorre $15\ \text{km}$ en $30\ \text{min}=0.5\ \text{h}$.
			\[
			v=\frac{d}{t}=\frac{15}{0.5}=30\ \text{km/h}
			\]
			
			\item \textbf{En m/s}
			
			Convertimos:
			\[
			15\ \text{km}=15000\ \text{m},
			\qquad
			30\ \text{min}=1800\ \text{s}
			\]
			\[
			v=\frac{15000}{1800}=\frac{25}{3}\approx 8.33\ \text{m/s}
			\]
		\end{itemize}
		
		\textbf{Respuesta:} $30\ \text{km/h}$ y aproximadamente $8.33\ \text{m/s}$.
		
		\item \textbf{Agua requerida por el cultivo}
		
		Cada planta requiere $2.5\ \text{L}$ diarios y hay 350 plantas:
		\[
		2.5\cdot350=875\ \text{L}
		\]
		
		Ahora convertimos a metros cúbicos:
		\[
		875\ \text{L}=\frac{875}{1000}=0.875\ \text{m}^3
		\]
		
		\textbf{Respuesta:}
		\begin{itemize}
			\item Se necesitan $875\ \text{L}$ diarios.
			\item Eso equivale a $0.875\ \text{m}^3$.
		\end{itemize}
	\end{enumerate}
	
	% =========================
	% SOLUCIONES IV
	% =========================
	\section*{IV. Soluciones: Expresiones algebraicas}
	
	\begin{enumerate}
		\item Simplifique las siguientes expresiones algebraicas:
		
		\begin{enumerate}[label=\alph*)]
			\item
			\[
			120mn - 54mn + 100m - 30mn - 28m
			\]
			Agrupando términos semejantes:
			\[
			(120-54-30)mn + (100-28)m = 36mn + 72m
			\]
			\textbf{Respuesta:} \(\boxed{36mn + 72m}\)
			
			\item
			\[
			35pq + 30 + 78pq - 60pq - 13
			\]
			\[
			(35+78-60)pq + (30-13)=53pq+17
			\]
			\textbf{Respuesta:} \(\boxed{53pq + 17}\)
			
			\item
			\[
			14 + [12xy - (3x^2y^2 - 4xy) + 8xy + 10]
			\]
			Distribuyendo el signo menos:
			\[
			14 + [12xy - 3x^2y^2 + 4xy + 8xy + 10]
			\]
			\[
			14 + [-3x^2y^2 + 24xy + 10]
			\]
			\[
			-3x^2y^2 + 24xy + 24
			\]
			\textbf{Respuesta:} \(\boxed{-3x^2y^2 + 24xy + 24}\)
			
			\item
			\[
			-10t + [ -(-2t + 3) + 12] - (14 + 2t + 11t)
			\]
			Primero:
			\[
			-(-2t+3)=2t-3
			\]
			entonces:
			\[
			-10t + (2t-3+12) - (14+13t)
			\]
			\[
			-10t + (2t+9) -14 -13t
			\]
			\[
			(-10t+2t-13t)+(9-14)=-21t-5
			\]
			\textbf{Respuesta:} \(\boxed{-21t - 5}\)
			
			\item
			\[
			x^2yz - (18xy^2z + 10xyz^2)
			\]
			Distribuyendo el signo menos:
			\[
			x^2yz - 18xy^2z - 10xyz^2
			\]
			No hay términos semejantes.
			
			\textbf{Respuesta:} \(\boxed{x^2yz - 18xy^2z - 10xyz^2}\)
			
			\item
			\[
			11a^2b - 32a^2b + 14ab^2 - 3a^2b
			\]
			\[
			(11-32-3)a^2b + 14ab^2 = -24a^2b + 14ab^2
			\]
			\textbf{Respuesta:} \(\boxed{-24a^2b + 14ab^2}\)
			
			\item
			\[
			10x^2yz - 12x^2yz - 8yx^2z + 42yx^2z
			\]
			Observamos que:
			\[
			yx^2z = x^2yz
			\]
			Entonces:
			\[
			(10-12-8+42)x^2yz = 32x^2yz
			\]
			\textbf{Respuesta:} \(\boxed{32x^2yz}\)
			
			\item
			\[
			54rt^2 - (12rt^2 - 5rt) + (21rt^2 - 2rt)
			\]
			Distribuyendo:
			\[
			54rt^2 - 12rt^2 + 5rt + 21rt^2 - 2rt
			]
			\[
			(54-12+21)rt^2 + (5-2)rt = 63rt^2 + 3rt
			\]
			\textbf{Respuesta:} \(\boxed{63rt^2 + 3rt}\)
		\end{enumerate}
	\end{enumerate}
	
	\newpage
	
	% =========================
	% SOLUCIONES V
	% =========================
	\section*{V. Soluciones: Problemas aplicados}
	
	\begin{enumerate}[itemsep=0.9cm]
		
		\item \textbf{Costo total}
		
		\[
		A = 3x^2 + 5xy - 2xy^2
		\]
		\[
		B = 10xy + 7x^2 - 3y^2
		\]
		\[
		C = 4y^2 - 8xy + 10
		\]
		
		Entonces:
		\[
		A+B+C=(3x^2 + 5xy - 2xy^2)+(10xy + 7x^2 - 3y^2)+(4y^2 - 8xy + 10)
		\]
		
		Agrupamos términos semejantes:
		\[
		(3x^2+7x^2)+(5xy+10xy-8xy)-2xy^2+(-3y^2+4y^2)+10
		\]
		\[
		10x^2 + 7xy - 2xy^2 + y^2 + 10
		\]
		
		\textbf{Respuesta:} \(\boxed{10x^2 + 7xy - 2xy^2 + y^2 + 10}\)
		
		\item \textbf{Crecimientos vegetales}
		
		\[
		A = 2x^3 - x^2 + 4x - 5
		\]
		\[
		B = x^3 - 8x - 2
		\]
		\[
		C = x^2 - 4x - 2
		\]
		
		Se pide:
		\[
		4A - [2B + C - (3B + 2A)] +10B
		\]
		
		Primero simplificamos el corchete:
		\[
		2B + C - (3B + 2A)=2B+C-3B-2A = -B + C - 2A
		\]
		
		Entonces:
		\[
		4A -[-B + C -2A] + 10B = 4A + B - C + 2A + 10B
		\]
		\[
		=6A + 11B - C
		\]
		
		Calculamos:
		\[
		6A = 6(2x^3 - x^2 + 4x - 5)=12x^3 - 6x^2 + 24x - 30
		\]
		\[
		11B = 11(x^3 - 8x -2)=11x^3 - 88x - 22
		\]
		\[
		-C = -(x^2 - 4x -2)= -x^2 + 4x + 2
		\]
		
		Sumando:
		\[
		(12x^3+11x^3)+(-6x^2-x^2)+(24x-88x+4x)+(-30-22+2)
		\]
		\[
		23x^3 - 7x^2 - 60x - 50
		\]
		
		\textbf{Respuesta:} \(\boxed{23x^3 - 7x^2 - 60x - 50}\)
		
		\item \textbf{Fertilización en un cultivo}
		
		Base:
		\[
		2n+15
		\]
		Refuerzo:
		\[
		3n-10
		\]
		
		Cantidad total:
		\[
		(2n+15)+(3n-10)=5n+5
		\]
		
		\textbf{Respuesta:} \(\boxed{5n+5}\)
		
		\item \textbf{Nutrientes}
		
		\[
		S = 45x^2 - 12xy + 30y^2
		\]
		\[
		F = -7x^2 + 24xy - 3y^2
		\]
		\[
		A = 30x^2 - 10xy + y^2
		\]
		
		\begin{itemize}
			\item \textbf{Con abono}
			
			\[
			S+F+A
			\]
			\[
			(45x^2 - 12xy + 30y^2)+(-7x^2 + 24xy - 3y^2)+(30x^2 - 10xy + y^2)
			\]
			\[
			(45-7+30)x^2+(-12+24-10)xy+(30-3+1)y^2
			\]
			\[
			68x^2 + 2xy + 28y^2
			\]
			
			\textbf{Respuesta con abono:} \(\boxed{68x^2 + 2xy + 28y^2}\)
			
			\item \textbf{Sin abono}
			
			\[
			S+F=(45x^2 - 12xy + 30y^2)+(-7x^2 + 24xy - 3y^2)
			\]
			\[
			(45-7)x^2+(-12+24)xy+(30-3)y^2
			\]
			\[
			38x^2 + 12xy + 27y^2
			\]
			
			\textbf{Respuesta sin abono:} \(\boxed{38x^2 + 12xy + 27y^2}\)
		\end{itemize}
		
		\item \textbf{Pesticidas}
		
		Disponible:
		\[
		50a + 20b + 30
		\]
		
		Aplicación total:
		\[
		(5a - 35b + 42)+(32a + 10b - 22)
		\]
		\[
		37a - 25b + 20
		\]
		
		Entonces, lo que queda es:
		\[
		(50a + 20b + 30) - (37a - 25b + 20)
		\]
		\[
		50a + 20b + 30 - 37a + 25b - 20
		\]
		\[
		13a + 45b + 10
		\]
		
		\textbf{Respuesta:} \(\boxed{13a + 45b + 10}\)
		
	\end{enumerate}
	
\end{document}